\section{Monochromatic Light}
\label{sec:geometry}

At radio frequencies, it is both possible to directly sample the electric field, and sufficient to employ a classical description of electromagnetic radiation \citep{Maxwell1865,bw80}.
%
At higher frequencies in the electromagnetic spectrum, polarimeters incorporate devices that count photons, such as charge-coupled devices and photomultiplier tubes.
%
For a comprehensive introduction to the quantum mechanical treatment of the polarization of light, please see
the excellent review by \citet{2021AdOP...13....1G}.

In both classical and quantum approaches to the topic, it is useful to develop some intuition for the geometry of polarization by starting with the nonphysical, idealized case of a purely monochromatic, plane-propagating, transverse electromagnetic wave.
%
For such a wave, there exist two linearly-independent solutions to Maxwell's equations, representing two oppositely polarized waves.  
%
Therefore, to fully describe the vector state of any observed signal, radio receiver systems are designed with a pair of receptors that are ideally sensitive to orthogonal senses of polarization.
%
Typically, the receptors are either circularly polarized (e.g.,  left and right circularly polarized)
or linearly polarized (e.g.,  horizontally and vertically polarized).

The following discussion is based on a Cartesian coordinate system in which the electromagnetic wave propagates in the positive $z$ direction and the electric field is completely described by two orthogonal linearly-polarized components in the $x$-$y$ plane, which is perpendicular to the direction of wave propagation,
\begin{equation}
	\Jcol{\epsilon}(t) =
	\left( 
	\begin{array}{c}
		\epsilon_x(t) \\
		\epsilon_y(t)
	\end{array}
	\right) =
	\left( 
	\begin{array}{c}
		a_x \cos (2\pi \nu t + \phi_x) \\
		a_y \cos (2\pi \nu t + \phi_y)
	\end{array}
	\right). \label{eqn:real_valued_wave}
\end{equation}
Here, $\epsilon_x$ and $\epsilon_y$ are the real-valued
projections of $\Jcol{\epsilon}$ onto the $x$ and $y$ axes.
%
In the special case of monochromatic light, $\nu$ is the constant frequency of the wave, $a_x$ and $a_y$ are the constant amplitudes of the orthogonal wave components, and $\phi_x$ and $\phi_y$ are the phases of these components at $t=0$.

On each cycle of the wave, the electric field vector traces an ellipse in the $x$--$y$ plane, as shown in \Fig{polarization_ellipse}.
%
%
\begin{figure}
	\centerline{\includegraphics[width=0.7\linewidth]{polarization_ellipse.pdf}}
\caption{
A left-handed ellipse is traced by the electric field vector on each cycle of a monochromatic wave travelling in the positive $z$ direction, as observed when looking toward the source.
%
The position angle, $-\pi/2 \le \psi \le \pi/2$, describes the orientation of the semi-major axis of the ellipse with respect to the $x$ axis; it is positive when the $x'$ and $y'$ axes are rotated in a counter-clockwise direction with respect to the $x$ and $y$ axes.
%
The tangent of the ellipticity angle, $-\pi/4 \le \chi \le \pi/4$, is equal to the ratio between the semi-minor axis and the semi-major axis; it is positive for a left-hand polarized wave and negative for a right-hand polarized wave.
%
The size of the ellipse is defined by $r$, which is equal to the length of the hypotenuse of the right triangle that contains $\chi$.}
\label{fig:polarization_ellipse}
\end{figure}
%
This figure depicts two angles that describe the geometry of this ellipse: the position angle $\psi$ and the ellipticity angle $\chi$.
%
\cite{sto52} first demonstrated the utility of expressing the components of the electric field vector in 
terms of these angles, starting with
%
an intermediate reference frame that is rotated about the 
$z$-axis by the position angle $\psi$, such that
\begin{eqnarray}
	x' &=& x \cos\psi + y\sin\psi \label{eqn:psi_rotation} \\
    y' &=& -x \sin\psi + y\cos\psi.
\end{eqnarray}
%
In this reference frame, which is depicted with dashed lines
and open arrow heads in \Fig{polarization_ellipse},
the electric field is expressed as in equation~(2) of \cite{sto52}
%
\begin{equation}
	\Jcol{\epsilon}'(t) =
	\left( 
	\begin{array}{c}
		\epsilon_x'(t) \\
		\epsilon_y'(t)
	\end{array}
	\right) =
	\left( 
	\begin{array}{c}
		r \cos\chi \sin (2\pi \nu t + \phi) \\
		r \sin\chi \cos (2\pi \nu t + \phi)
	\end{array}
	\right),
	\label{eqn:stokes_ellipse}
\end{equation}
where $r$ defines the size of the ellipse and $\phi$ defines the position of the electric field vector at $t=0$.

When $\chi = 0$, $\epsilon_y' = 0$ and the electric field vector oscillates along a line defined by the $x'$ axis, which
is oriented with respect to the $x$ axis by $\psi$.
%
When $\chi > 0$, 
%the phase of $\epsilon_y'$ leads that of $\epsilon_x'$ (because the phase of $\cos(2\pi\nu t)$ leads the phase of $\sin(2\pi\nu t)$ by $\pi/2$) and 
the electric field travels in a clockwise direction about the ellipse, as seen by an observer looking toward the source, which is defined as a left-hand polarized wave \citep{ieee145}.
%
% At the top of page 34 of IEEE Std 145-2013:
%
% "NOTE — When the plane of polarization is viewed from a specified side, if the extremity of the field vector rotates clockwise [counterclockwise] the sense is right-handed [left-handed]. For a plane wave, the plane of  polarization is viewed looking in the direction of propagation."
%
When $\chi < 0$, the counter-clockwise traverse of the electric field vector 
is defined as right-hand polarization.
%
In the special cases of $\chi = \pm \pi/4$, the amplitudes of $\epsilon_x'$ and $\epsilon_y'$ are equal and the electric field traces a circle in the $x$--$y$ plane.
%
In general, polarization state depends on both the amplitudes of the electric field components and the relative phase between them.  

To move beyond the nonphysical case of perfectly monochromatic waves, for which amplitude and phase are constant as a function of time, it is necessary to characterize random fluctuations of both amplitude and phase using statistical averages.  This is facilitated by adopting the complex-valued analytic representation of a signal, which directly encodes its instantaneous amplitude and phase (see \App{analytic}).
%
% the imaginary part of the analytic signal is derived from the real part and contains no new information; furthermore, the real-part of the analytic signal is simply the original real-valued signal.
%
%
After replacing the real-valued $\epsilon_x(t)$ and $\epsilon_y(t)$ with their associated analytic signals, $e_x(t)$ and $e_y(t)$, the monochromatic wave described by \eqns{real_valued_wave}{through}{stokes_ellipse} is given by the real part of its analytic representation,
\begin{equation}
\label{eqn:analytic_field}
\Jcol{e}(t) =
	\left( 
	\begin{array}{c}
		e_x(t) \\
		e_y(t)
	\end{array}
	\right) = \Jcol{e}_0 \exp(\Ci [2\pi \nu t + \phi - \pi/2]).
\end{equation}
%
Here, the polarization state of the wave is completely described by the complex-valued two-dimensional Jones vector \eqnp{monochromatic_analytic},
\begin{equation}
\label{eqn:Jones_vector}
\Jcol{e}_0=r\left( \begin{array}{c}
    \cos\psi\cos\chi - \Ci\sin\psi\sin\chi \\
    \sin\psi\cos\chi + \Ci\cos\psi\sin\chi
  \end{array} \right).
\end{equation}
%
An interactive notebook\footnote{\nburl} demonstrates the relation between $\Jcol{\epsilon}(t)$ and $\Jcol{e}(t)$ by depicting the ellipses inscribed by these functions as $r$, $\psi$ and $\chi$ are varied.
